\section{Methods}

This section distinguishes implemented infrastructure from planned control
methods.

\subsection{Two-dimensional vehicle model}

\implemented{} The planar plant uses a low-speed kinematic bicycle for the
tractor and a no-lateral-slip single-axle trailer. Its independent state is
$[x_t,y_t,\psi_t,\psi_r,\delta,v]$. With wheelbase $L_t$, trailer length $L_r$,
signed longitudinal and lateral hitch offsets $a$ and $b$, tractor yaw rate
$\omega=v\tan(\delta)/L_t$, and articulation
$\alpha=\psi_t-\psi_r$, the trailer equation is
\begin{equation}
 \dot{\psi}_r = \frac{(v-b\omega)\sin\alpha+a\omega\cos\alpha}{L_r}.
\end{equation}
Together with $\dot{x}_t=v\cos\psi_t$, $\dot{y}_t=v\sin\psi_t$,
$\dot{\psi}_t=\omega$, $\dot{\delta}=u_\delta$, and bounded acceleration
$\dot v=u_v$, this reduces to the frozen on-axle model when $a=b=0$.

The simulator applies command supervision at $10\,\mathrm{Hz}$ and advances
five fixed $0.02\,\mathrm{s}$ classical fourth-order Runge--Kutta substeps.
Hitch and trailer-axle positions are reconstructed from rigid geometry rather
than independently integrated. Tests cover analytic derivatives, straight and
turning motion in both speed directions, articulation, angle wrapping, limits,
rigid reconstruction, deterministic replay, timestep convergence, and direct
low-speed trajectory comparison with MuJoCo.

\subsection{Reference trajectories and observations}

\implemented{} Reference paths are non-degenerate polylines ordered from the
dock target toward the start so a reverse maneuver moves toward decreasing arc
length while the nominal vehicle heading remains outward. Orthogonal projection
provides signed lateral error and progress; preview points are sampled toward
the goal at $0.5$, $1.0$, and $2.0\,\mathrm{m}$.

The bounded 21-element observation contains lateral error, sine/cosine tractor
and trailer heading errors, articulation and articulation rate, steering state,
speed, remaining path distance, docking-pose error, and the local coordinates
of the three preview points. Values are normalized and clipped to $[-1,1]$,
with a stable name assigned to every element.

\subsection{Gymnasium docking task}

\implemented{} The registered environment advances one $0.10\,\mathrm{s}$
plant command per step. Its primary action is normalized steering rate at the
fixed reverse-speed target. A separate two-action teleoperation mode adds speed
authority for manual debugging without changing the research action contract.
The reward combines progress with penalties for time, lateral and heading
error, articulation, and steering rate, plus explicit terminal terms. Every
component is returned in the diagnostic information and versioned in a strict
task configuration. Success, collision, jackknife, workspace exit, and timeout
are evaluated independently of reward.

The compact registered task remains the controller-comparison baseline. The
full-size interactive task is selected explicitly by separate vehicle, dock,
and task configurations. In that demonstration, the colored active loading
zone has no longitudinal side-wall collision polygons; complete-trailer
containment still uses the configured zone width. The rear boundary and
building façade outside the active door remain collidable. Its hard jackknife
event occurs at $60^\circ$, while the compact baseline terminates at
$45^\circ$.

\subsection{NMPC baseline}

\planned{} NMPC will predict the shared articulated-vehicle model and minimize
tracking, articulation, steering effort, and steering-rate variation while
enforcing dynamics, steering, steering-rate, and articulation constraints. The
primary comparison will retain fixed reverse speed and steering-rate authority.
Horizon, discretization, weights, warm start, solver tolerances, and numerical
failure handling will be recorded in versioned experiment configurations.

\subsection{PPO controller}

\planned{} PPO will use the same steering-rate authority and fixed reverse
speed. Training is expected to proceed from straight reverse stabilization to
curved tracking and randomized docking. Reward components, termination logic,
network architecture, training budget, and curriculum transitions will be
documented after implementation. A global-coordinate policy and removal of
preview information are candidate observation ablations, not established
experiments.

\subsection{Safety supervision}

\implemented{} Raw and supervised commands are represented separately so
interventions are observable. Steering-rate, steering-angle, speed, and
acceleration limits are enforced through the shared configuration. Chassis
polygons are checked against margin-adjusted dock walls, and hard collision and
articulation events are independent of learned reward.
